\chapter{Introduction}

\section{Motivation}
\small{
Robotic systems are increasingly deployed across industrial, service, medical, and consumer domains. According to the International Federation of Robotics, annual installations of industrial robots increased from approximately 221,000 units in 2014 to 542,000 units in 2024, indicating that global robot demand in factories has roughly doubled over the last decade\cite{ifr2025robots}. In the same period, the global operational stock of industrial robots increased from approximately 1.47 million units to 4.66 million units\cite{ifr2025robots}. This development shows that robots are becoming a more widespread component of modern production and service environments.

At the same time, robotic systems are increasingly influenced by advances in artificial intelligence. The Stanford AI Index reports that 78\% of organizations used AI in 2024, compared with 55\% in the previous year \cite{stanford2025aiindex}.In addition, global private investment in generative AI reached USD 33.9 billion in 2024. The International Federation of Robotics also identifies physical, analytic, and generative AI as key technological trends in robotics. These developments indicate that future robotic systems will increasingly combine physical operation with adaptive, data-driven, and software-defined capabilities\cite{ifr2025robots}.

This combination introduces new challenges for safety and risk assessment. Conventional risk assessment methods are typically performed before deployment and are based on assumptions about the system design, intended use, operating environment, and known hazards\cite{iso12100}. While such methods remain essential, they are limited when the system behaviour or operating environment changes after deployment. This limitation becomes particularly relevant for autonomous systems that operate in dynamic environments, receive over-the-air software updates, or use adaptive AI-based components.

As a result, there is a growing need for dynamic risk assessment mechanisms that can incorporate information about the current operational context. Such mechanisms require structured representations of the environment, including relevant objects, spatial relationships, and risk-related information. This motivates the development of a spatial risk-annotated scene graph, which can provide environmental context and situational awareness to support risk assessment during system operation.

}
\section{Relevance}

\section{Goal of the Thesis}